\documentclass[11pt,a4paper]{article}
\usepackage{fontspec}
\setmainfont{Noto Serif}
\setsansfont{Noto Sans}
\setmonofont{Noto Sans Mono}[Scale=0.84]
\usepackage[english]{babel}
\usepackage{amsmath,amssymb,mathtools}
\usepackage[margin=2.25cm,headheight=15pt]{geometry}
\usepackage{booktabs,longtable,tabularx,array,graphicx}
\usepackage{xcolor,hyperref,xurl,fancyhdr,enumitem}
\definecolor{ink}{HTML}{183449}
\definecolor{teal}{HTML}{187C80}
\definecolor{pale}{HTML}{E8F2F2}
\hypersetup{colorlinks=true,linkcolor=ink,urlcolor=teal,pdftitle={Week 2: HAPT Data Foundations},pdfauthor={ECDR Research}}
\urlstyle{same}
\graphicspath{{output/week2_foundations_20261103/}{../output/week2_foundations_20261103/}}
\setlength{\parindent}{0pt}\setlength{\parskip}{5pt}\setlength{\emergencystretch}{3em}
\setlength{\tabcolsep}{4pt}\renewcommand{\arraystretch}{1.13}
\newcolumntype{Y}{>{\raggedright\arraybackslash}X}
\newcolumntype{P}[1]{>{\raggedright\arraybackslash}p{#1}}
\newcommand{\R}{\mathbb R}\newcommand{\N}{\mathbb N}
\newcommand{\norm}[1]{\left\lVert#1\right\rVert}\newcommand{\abs}[1]{\left\lvert#1\right\rvert}
\newcommand{\code}[1]{{\small\ttfamily\nolinkurl{#1}}}
\newcommand{\ind}{\mathbf 1}\newcommand{\Sset}{\mathcal S}
\pagestyle{fancy}\fancyhf{}\fancyhead[L]{\small\textcolor{ink}{ECDR | Week 2 data foundations}}\fancyhead[R]{\small HAPT and frozen reference}\fancyfoot[L]{\small 3 November 2026 | English edition}\fancyfoot[R]{\thepage}
\setlist{nosep,leftmargin=1.5em}\raggedbottom
\begin{document}
\begin{titlepage}
\vspace*{1.2cm}{\Large\color{teal}DATA FOUNDATIONS AND AUDIT\par}\vspace{.5cm}
{\huge\bfseries\color{ink}HAPT Data Foundations\par}\vspace{.4cm}
{\Large Physical data contract, dependence, support, preprocessing, and decision evidence\par}\vspace{1cm}
\begin{tabularx}{\linewidth}{@{}P{3.1cm}Y@{}}\toprule
Dataset & Human Activity and Postural Transition (HAPT): 30 participants, 61 recordings, six inertial channels, 50 Hz.\\
Window tensor & 17,453 windows with shape $6\times128$, stride 64; float32 and finite.\\
Reconstruction & All rebuilt sensor values and all 10 manifest fields match the frozen window artifact exactly.\\
Reference & Week 1 selected K=12 SmallCNN1D, seed 43/epoch 8, 283,692 trainable parameters; this report describes its data contract and evidence boundary.\\
Key open item & Historical majority-support counts do not reproduce the current source predicate. The cause remains unproven.\\
\bottomrule\end{tabularx}
\vfill
This is an independent Week 2 data study. Its objective is to specify and audit what enters the frozen reference pipeline: sampled physical channels, window coordinates, annotation policy, subject partitions, and train-derived normalization. It asks whether those data objects can be reconstructed from the identified source and whether the resulting handoff is unambiguous. The method combines source inspection with a read-only reconstruction audit against frozen tensors, manifests, split metadata, and normalization values. Conclusions are limited to that data contract; network architecture, prediction accuracy, and certification mathematics belong to the separate Week 2 CNN study.
\end{titlepage}
\tableofcontents\clearpage

\section{Study contract and evidence levels}
\textbf{Abstract.} This study reconstructs the HAPT data path used by the frozen Week 1 reference. It specifies six physical sensor channels, the 128-sample/64-stride window rule, sample-index label policy, subject-disjoint partitions, supervised denominators, and train-only normalization. A read-only audit reproduces 17,453 windows and every manifest field; all 17,392 eligible aligned overlap pairs match exactly. The study finds that the current local pipeline is reproducible while upstream annotation index origin and historical support-file provenance remain unresolved. Transition classes have much lower support than basic activities, and overlapping windows are dependent observations. These results establish a bounded local data contract, not model quality or certification.

The source revision is \code{633dc05a80511af4d4580cd985edd494cdc373ff} on \code{main}. The worktree was clean before this Week 2 revision. Local report and audit additions are recorded by hash in the Week 2 final verification receipt. This report is the current standalone Week 2 authority; the older combined Weeks 1--3 document remains a historical synthesis.

The study objective is to make explicit what data the frozen Week 1 reference receives, how each window and label is built, how examples are divided and normalized, and which claims can be supported by the resulting artifacts. The answer is bounded to the canonical current source and the existing run. It is not a new training study and does not revise the selected checkpoint, labels, splits or preprocessing.

The executed read-only audit is \code{scripts/week2_data_audit.py}. It rebuilt windows from 122 raw sensor files and \code{labels.txt}, checked the reconstructed tensor and manifest against the frozen artifacts, recomputed subject-disjoint split and supervised support, and independently recomputed the train-only normalization. The full machine-readable receipt is \code{output/week2_foundations_20261103/data_audit.json}; row-level adjacent overlap checks are in \code{data_aligned_overlap_pairs.csv}; the historical-support comparison is \code{data_support_comparison.json}; interpreted results are in \code{data_findings.md}. The audit records its environment, timestamp and input hashes. The current model behavior audit is separate and recorded in \code{model_audit.json} and \code{model_findings.md}.

Four evidence labels are useful throughout. \emph{Source contract} means behavior read from an identified implementation or locked source-of-truth file. \emph{Verified artifact} means a check was actually run against identified files. \emph{Mathematical consequence} means a result follows under named definitions and assumptions. \emph{Open provenance} means the available files do not resolve a historical or semantic question. A checked checksum proves file identity, not the truth of an undocumented historical interpretation. A successful reconstruction proves the implementation reproduces its own frozen artifact, not that every upstream annotation convention is semantically correct.

The canonical Week 1 contract is \code{00_OVERVIEW/week1/SOURCE_OF_TRUTH.md}. The data-construction path includes \code{scripts/03_build_windows.py}, \code{ecdr/data.py}, \code{ecdr/preprocessing.py}, \code{scripts/ecdr_audit_support.py}, and \code{configs/backbone.yaml}. The pipeline's currently used annotation slice begins at the zero-based raw index used by its window builder. The UCI documentation describes the dataset but does not establish the index-origin choice used by this local builder. An independently executed one-sample offset changes 98 window policy/label assignments; this makes the convention consequential while leaving its intended semantic origin unresolved.

The dataset is distributed by the \href{https://archive.ics.uci.edu/dataset/341/smartphone+based+recognition+of+human+activities+and+postural+transitions}{UCI Machine Learning Repository}. That source describes the study and file families. For exact claims in this report, the checked out source, frozen manifests and audit receipt govern the local pipeline behavior. This distinction prevents a public dataset description from being used to silently fill a missing local provenance link.

\section{W2-01: Physical recordings and channel semantics}
HAPT contains body-worn inertial measurements collected from 30 subjects performing six basic activities and six postural transitions. The local data inventory has 61 recordings, each with one accelerometer file and one gyroscope file, for 122 sensor files in total. The raw sample rate is 50 Hz. Each raw time step supplies six row-aligned values: three acceleration axes and three angular-velocity axes.

For recording $r$, denote its signal by
\[
 X^{(r)}=(x^{(r)}_{t,c})_{t=0}^{T_r-1}{}_{,\,c=0}^{5}\in\R^{T_r\times6}.
\]
Row alignment is the local data contract; equality of file dimensions and row indices does not independently verify physical sensor-clock synchronization. The first three columns correspond to acceleration axes $x,y,z$ and the final three to gyroscope axes $x,y,z$. Acceleration is stored in multiples of gravitational acceleration $g$; gyroscope angular velocity is stored in radians per second (rad/s). The report preserves the dataset's acceleration unit and does not convert it to metres per second squared. This formula represents the raw temporal array; the model's window tensor uses channel-first convention, so a window has shape $C\times W=6\times128$. A batch has shape $B\times6\times128$, with batch size $B$.

Accelerometer and gyroscope components have different stored physical units and distributions. Their numerical magnitudes cannot be pooled into a single physical Euclidean distance without first defining unit conversion, scaling and intended semantics. The Week 2 raw-channel table therefore reports each channel separately. The later network normalization makes each input coordinate dimensionless relative to training-set statistics, but that standardized coordinate system does not retroactively change the physical units of the raw source measurements.

The HAPT distribution also provides a 561-dimensional engineered-feature representation derived from sensor windows. Those feature dimensions are not 561 additional raw files and are not concatenated to the six-channel signal in the frozen model input. The model receives the six raw inertial channels after windowing and channelwise normalization. Counting the engineered files as extra channels would misstate the reference function and its input contract.

The sensor channels describe measured motion; the activity labels are annotations aligned to sample indices. A label is not another sensor measurement, and a transition label does not imply that the six numeric channels encode a unique event boundary. Label alignment and sensor-time indexing are separate pieces of provenance. This distinction becomes important where a one-sample change to the alignment convention alters derived labels even though the raw sensor values remain unchanged.

\section{W2-02: Window construction, duration, and overlap}
The builder takes fixed-length windows of $W=128$ samples with stride $s=64$. For a recording of length $T$ and the convention that only complete windows are emitted, the window count is
\[
 n(T;W,s)=\begin{cases}0,&T<W,\\1+\left\lfloor\dfrac{T-W}{s}\right\rfloor,&T\ge W.\end{cases}
\]
This follows because the start indices are $0,s,2s,\ldots,qs$ with $qs+W\le T$; the largest admissible integer $q$ is $\lfloor(T-W)/s\rfloor$. Summing this count over the 61 recordings produces 17,453 windows. No window crosses from one recording into another: the start index restarts within each recording, and the unused trailing segment of each recording is not joined to the next file.

A window is
\[
 X_i[c,u]=x^{(r)}_{a_i+u,c},\quad c=0,\ldots,5,\quad u=0,\ldots,127,
\]
where $r$ identifies its recording and $a_i$ its start sample within that recording. At 50 samples per second, 128 sample intervals correspond to $128/50=2.56$ seconds under the dataset window-duration convention. The elapsed time from the first sample timestamp to the last is instead $(128-1)/50=2.54$ seconds because 127 inter-sample gaps separate 128 endpoints. Both numbers are correct for their stated interpretation; a duration claim must say whether it counts sample slots or endpoint separation. A stride of 64 samples is $64/50=1.28$ seconds, half the 2.56-second window duration.

For adjacent windows of the same recording, $a_{i+1}=a_i+64$. The first window covers indices $a_i$ through $a_i+127$, while the next covers $a_i+64$ through $a_i+191$. Their shared raw indices are $a_i+64$ through $a_i+127$, exactly 64 samples. In tensor-local coordinates this identity is
\[
 X_i[c,64+q]=X_{i+1}[c,q],\qquad q=0,\ldots,63.
\]
Equivalently, write $X_{i+1}=[X_i[:,64:128],U_{i+1}]$, where $U_{i+1}\in\R^{6\times64}$ is the newly appended raw block. This shift--append identity specifies the raw input only; it does not claim reusable network activations or unchanged logits.

The audit checked every eligible pair: 17,392 aligned adjacent pairs, all 64-by-6 shared values exactly equal in raw storage and, after applying the frozen normalization, exactly equal in normalized storage. This is a deterministic consequence of cutting overlapping windows from the same recording and applying the same fixed coordinatewise transform.

The incoming portion of a stride-aligned next window has 64 samples per channel, hence $64\times6=384$ new scalar sensor values. In the frozen float32 tensor, this is $384\times4=1,536$ bytes, or 1.5 KiB. In float64 it is $384\times8=3,072$ bytes, or 3 KiB; the representation must accompany any memory estimate. The prior 64-sample overlap need not be reacquired if the system owns a valid raw ring buffer, but the complete six-channel float32 window contains 768 values (3,072 bytes). This byte count describes a tensor payload only; it excludes metadata, allocator alignment, copies and other process memory. The model forward may process the complete window even if the acquisition layer reuses prior samples.

Exact coordinate overlap does not mean the two windows are numerically close at equal local coordinates. For $u=0,\ldots,127$, local coordinate $u$ in $X_i$ represents raw time $a_i+u$, while coordinate $u$ in $X_{i+1}$ represents $a_i+64+u$. Equal-local-coordinate comparison therefore contrasts different physical times. For a simple scalar ramp $x_t=t$, these two windows can share 64 exact physical samples at offset coordinates, yet $X_i[u]-X_{i+1}[u]=-64$ wherever both values exist at equal local index. The window tensors are not close in that coordinate comparison even though overlap identity holds exactly.

The audit's descriptive mean absolute difference at equal local coordinates was computed independently by channel. In raw stored units, values were 0.15258, 0.11625 and 0.10393 for accelerometer $x,y,z$, and 0.43151, 0.34394 and 0.25361 for gyroscope $x,y,z$. In normalized dimensionless coordinates the corresponding values were 0.37602, 0.28552, 0.29283 and 0.96399, 0.81642, 0.85550. These are temporal-shift comparisons over aligned pairs, not distances between the two overlapping physical regions. A pooled six-channel norm is intentionally omitted because acceleration and angular velocity have distinct physical units before standardization, and because a single norm would encode an additional scaling choice.

The full audit inventory also records per-channel medians and 95th percentiles for this equal-local-coordinate statistic, as well as each pair's values, in the CSV and JSON. These are descriptive features of this recording set and windowing rule. They do not by themselves define an ECDR candidate region or prove a tolerance that preserves decisions.

\section{W2-03: Label policy and annotation alignment}
Let $y^{(r)}_t$ denote the annotation class at raw sample $t$ when one is present. For a window, let $n_k$ be the number of its 128 raw annotation indices assigned to class $k$, for $k\in\{1,\ldots,12\}$. Let $n_0$ count unannotated or excluded samples, and let $n_{max}=\max_{k=1}^{12}n_k$. The current window policy distinguishes four cases.
\begin{description}
\item[Pure.] Exactly one activity class covers the full window: $n_k=128$ for one $k$. It receives that class label and policy tag \code{pure}.
\item[Majority.] The dominant annotated class covers at least 60 percent of the full window: $n_{max}/128\ge0.6$. Because $0.6\times128=76.8$, an integer count must have at least 77 samples. The policy is \code{majority}, with the winning class ID.
\item[Ambiguous.] The window contains annotation but no class meets the majority threshold; its policy tag is \code{ambiguous} and it is excluded from supervised K=12 training/evaluation.
\item[Unlabeled.] No eligible activity annotation is available for the window under the builder's convention. The row is retained in the full stream with the \code{unlabeled} policy and does not enter supervised metrics.
\end{description}
The exact implementation controls tie behavior. A dominant-count rule that returns one class under a tie is not equivalent to a policy that rejects ties. For the accepted-majority set, any tie at the maximum would still need to satisfy the 77-of-128 threshold; two tied classes cannot each have 77 samples, so an accepted tie is arithmetically impossible. However, details of how the builder stores a dominant label for ambiguous rows do not make those rows supervised. Their inclusion is governed by the explicit policy tag as well as label value.

The label uses a fixed denominator of 128 raw window positions for the 60 percent threshold. It is not a percentage of only annotated positions. For example, a window with 77 samples of class 4 and 51 samples from another annotation satisfies the majority rule for class 4. A window with 76 samples of class 4 and 52 samples from other classes does not; even though class 4 is dominant, it falls below the integer threshold and remains ambiguous. If unannotated positions occur, the full-window denominator remains the governing contract where implemented; one must not silently change to $n_{max}/(128-n_0)$.

The current frozen manifest has 10 fields per row and 17,453 rows. Reconstruction compared all ten fields, not only labels: recording and subject identity, window index, start/end coordinates, policy and class summary fields are part of the row-level trace. Split membership is not a manifest column; it is reconstructed and checked separately by joining on recording and window identity. In addition, the reconstructed raw tensor matched the frozen float32 tensor value for value. Passing a total count check alone would not detect a wrong row order, wrong label, changed subject assignment, or swapped start index; the full-row equality and split-join checks address these distinct failure modes.

Annotation indexing remains a provenance limitation. The builder treats each labels-file start/end range as inclusive when assigning its annotation across raw samples, then slices those labels with the same zero-based positions used to form each sensor window. That convention yields the current manifest. A diagnostic one-sample shift in label indexing changes 98 window policy or label assignments. This proves the derived artifact is sensitive to a one-sample convention. It does not prove that the shifted result is more correct, and the public README does not establish the original annotation index origin clearly enough to resolve the choice. The proper status is “current implementation reproduced; upstream semantic index origin OPEN.”

\section{W2-04: Splits, denominators, and class support}
Subjects, not individual windows, define the frozen split boundary. The 30 subject IDs are divided into 17 train, 4 validation and 9 test subjects. Recordings map wholly to one split, and the subject groups are pairwise disjoint. This reduces direct participant leakage between those three partitions. It does not make windows within a partition independent: adjacent windows share half their raw samples, multiple recordings can belong to one person, and subjects may share population-level characteristics.

There are two denominators because the sequential stream keeps every window while supervised classification uses only pure and accepted-majority rows with one of the 12 class labels. The split counts reconcile as follows.
\begin{center}\begin{tabular}{@{}lrrrr@{}}\toprule
Split & Full stream & Supervised & Ambiguous excluded & Unlabeled excluded\\\midrule
Train & 9,978 & 6,955 & 904 & 2,119\\
Validation & 2,422 & 1,751 & 205 & 466\\
Test & 5,053 & 3,572 & 479 & 1,002\\\midrule
Total & 17,453 & 12,278 & 1,588 & 3,587\\\bottomrule
\end{tabular}\end{center}
Each row satisfies full stream = supervised + ambiguous + unlabeled. Across the entire tensor, policy totals are 10,361 pure, 1,917 majority, 1,588 ambiguous and 3,587 unlabeled, summing to 17,453. Therefore the supervised count is pure plus majority: $10,361+1,917=12,278$. Unlabeled and ambiguous rows remain part of the sequential stream, but do not supply a single accepted target for supervised analysis.

The selected model's train/validation/test class supports are:
\begin{center}\small\begin{tabular}{@{}rP{4.6cm}rrr@{}}\toprule ID & Class & Train & Validation & Test\\\midrule
1 & WALKING & 1,077 & 246 & 531\\
2 & WALKING\_UPSTAIRS & 973 & 247 & 537\\
3 & WALKING\_DOWNSTAIRS & 894 & 231 & 492\\
4 & SITTING & 1,103 & 289 & 546\\
5 & STANDING & 1,206 & 312 & 597\\
6 & LAYING & 1,188 & 316 & 586\\
7 & STAND\_TO\_SIT & 79 & 19 & 45\\
8 & SIT\_TO\_STAND & 53 & 13 & 32\\
9 & SIT\_TO\_LIE & 90 & 26 & 51\\
10 & LIE\_TO\_SIT & 89 & 15 & 46\\
11 & STAND\_TO\_LIE & 119 & 18 & 68\\
12 & LIE\_TO\_STAND & 84 & 19 & 41\\\bottomrule\end{tabular}\end{center}
Test support for the six basic activities ranges from 492 to 597, while transition support ranges from 32 to 68. The complete three-split support is retained in the machine-readable audit. An aggregate of 17,453 windows can look large, but it does not remove the sparse evidence for some transitions. The Week 1 authority records the transition-support quality gate as HOLD; this Week 2 count audit does not change that decision.

\subsection{Historical majority-count discrepancy}
The current source function \code{ecdr.data.audit_12_class_support}, reached through \code{scripts/ecdr_audit_support.py}, selects rows whose policy is in $\{\texttt{pure},\texttt{majority}\}$ and whose class ID matches the requested class. On the frozen manifest it yields supervised totals 6,955, 1,751 and 3,572. The pure-only totals are 5,878, 1,490 and 2,993. The historical file \code{data/ecdr_12_class_support.json} records values called majority of 7,859, 1,956 and 4,051.

Those historical values exceed the current supervised totals by 904, 205 and 479. Those differences equal the current ambiguous counts in the corresponding splits. This numerical correspondence is an observed relation between two snapshots; it is not sufficient evidence that the historical generator included ambiguous rows, used a particular predicate, or intended a different label policy. The historical JSON has no run provenance or implementation record sufficient to decide how it was produced. It cannot be relabeled as current execution evidence, and its blocker field/name does not prove model acceptance. We preserve the mismatch as OPEN historical provenance rather than invent a cause. The exact current predicate and comparison are serialized in \code{data_support_comparison.json}.

\section{W2-05: Train-only normalization and coordinate systems}
For supervised training windows $\mathcal I_{tr}$, the frozen pipeline computes per-channel statistics over the window and time axes. With $N_{tr}=6,955$ supervised windows, define
\[
 \mu_c=\frac{1}{N_{tr}W}\sum_{i\in\mathcal I_{tr}}\sum_{t=0}^{127}X_i[c,t],\qquad
 \sigma_c=\sqrt{\frac{1}{N_{tr}W}\sum_{i\in\mathcal I_{tr}}\sum_{t=0}^{127}(X_i[c,t]-\mu_c)^2}.
\]
The audit used float64 accumulation, stored statistics as float32, and applies the configured floor $\sigma_c^*=\max(\sigma_c,10^{-6})$. The normalized value is
\[
 Z_i[c,t]=\frac{X_i[c,t]-\mu_c}{\sigma_c^*},
\]
cast to float32 for the model input. Validation and test use these same training-derived numbers. They do not estimate their own means or standard deviations.

\begin{center}\begin{tabular}{@{}lrr@{}}\toprule Channel & Training mean & Training standard deviation\\\midrule
Acceleration X & 0.8043609262 & 0.4057796597\\
Acceleration Y & 0.0411366001 & 0.4071546793\\
Acceleration Z & 0.0817700475 & 0.3549118936\\
Gyroscope X & 0.0091355676 & 0.4476351142\\
Gyroscope Y & $-0.0043975562$ & 0.4212706089\\
Gyroscope Z & $-0.0071832035$ & 0.2964415252\\\bottomrule
\end{tabular}\end{center}
All six recomputed statistics match the saved \code{normalization.json} and selected-checkpoint metadata exactly. The preprocessing identity is \code{aa8e1cc20829307187e0a78200f9957f7dd068b1dc26063c2720ad832bb2fd6e}. The audit confirms the identity and arrays; it does not claim the statistics are a universal physical calibration.

The train statistics are window-weighted, not unique-raw-time-weighted. If a raw sample occurs in $m_t$ overlapping supervised training windows, it contributes $m_t$ times to the sum above. For the 50-percent-overlap scheme, many interior samples occur in two adjacent windows, while boundary or excluded-window samples can occur zero or one time. This weighting follows the actual training loader and is part of the frozen preprocessing contract. Deduplicating shared raw indices would yield a different mean and variance and therefore a different model function.

Because normalization is a fixed shared affine map, exact shared raw overlap remains exact after normalization:
\[
 X_i[c,64+q]=X_{i+1}[c,q]\;\Longrightarrow\;
 Z_i[c,64+q]=Z_{i+1}[c,q].
\]
This implication requires both windows to use the same channel constants and same arithmetic. It holds for the frozen transform audited here. It would not generally hold if each window independently estimated its own mean and scale.

Normalization changes coordinate magnitudes but does not turn an arbitrary distance into a physically meaningful one. If a one-channel raw perturbation is $\delta x_c$, its normalized magnitude is $\delta z_c=\delta x_c/\sigma_c^*$. For example, using the stored acceleration-X standard deviation, a raw-coordinate change of 0.4057796597 stored units maps to approximately 1 normalized unit. The analogous one-standard-deviation raw change differs for every channel. This example illustrates units; it does not define an allowed perturbation radius or imply that such a change is behaviorally harmless.

A separate function in \code{ecdr/preprocessing.py} accumulates float64, returns a float64 transform, defaults to $\sigma_{min}=10^{-5}$, and hashes a JSON contract (including channel names/dimensions) to a truncated 16-character identifier. The frozen Week 1 model path uses float32 stored statistics, floor $10^{-6}$, float32 normalized tensors, and a byte-based full SHA-256 preprocessing ID. Even if the floor is inactive for all six observed channel scales, the numeric/data-type and identity contracts differ. Equal values on this particular dataset would not make the two functions interchangeable for every input. Their IDs must not be compared as if generated from the same serialization rule, and a model checkpoint bound to one preprocessing identity cannot silently be paired with the other.

\section{W2-06: Data-quality evidence and dependence}
This unit examines what the windowing and annotation records permit us to say about data quality. It does not infer prediction quality. The adjacent-window audit distinguishes physical overlap identity from equal-local-coordinate change. The exact 64-sample shared region is identical in every eligible pair; comparing the same local coordinate in successive windows instead compares measurements 64 samples apart in the recording. Its per-channel mean absolute difference is reported in W2-02. Because acceleration and angular velocity have different raw units, those channelwise results remain separate and are not pooled into a physical norm.

The audit also computes within-recording, channelwise raw-sample correlations at lags 64 and 128. Mean correlations are 0.43783 and 0.36346, respectively. At lag 64, acceleration X/Y/Z are 0.72095/0.82400/0.77693, while gyroscope X/Y/Z are 0.14840/0.06063/0.09604. These are descriptive statistics over the audited recordings. They are not an independence test, a population model, or uncertainty intervals. In particular, the 50-percent shared raw input between adjacent windows is direct evidence that window rows cannot be treated as independent samples.

There are 30 subjects partitioned into non-overlapping train, validation, and test groups of 17, 4, and 9. Keeping subjects disjoint prevents the same identified participant from appearing in multiple partitions, but it does not establish that participants are independent draws from a specified target population. Recordings within a person and adjacent windows within a recording retain temporal and subject structure. For this reason, the study reports both stream and supervised denominators and descriptive class counts; it does not turn window totals into independent participant counts or use them to claim population confidence.

The basic-activity classes have 492--597 supervised test windows each, whereas transition classes have 32--68. This support imbalance is directly visible in the manifest-derived counts and is consistent with the existing Week 1 G2 quality status of HOLD. The present audit documents that status and its count basis; it does not revise the gate or attribute the imbalance to a cause. Further claims about annotation quality would require an authorized reference or adjudication procedure that is not present in these artifacts.

\section{W2-07: Dataset provenance and reproducibility}
The audit establishes exact local reconstruction: 122 sensor files plus the label file reproduce all values in the frozen window tensor and all ten manifest fields. This is a strong reproducibility result for the identified local builder and frozen artifacts. It does not independently validate the upstream meaning of every sample coordinate or label. The builder uses its current integer-index convention; shifting the annotation coordinates by one sample changes 98 window policy or label assignments. The audit demonstrates sensitivity to the convention, while the semantic origin intended by the upstream annotation source remains OPEN.

A second unresolved record is the historical support snapshot \code{data/ecdr_12_class_support.json}. Its recorded majority totals are 7,859/1,956/4,051, whereas the current source predicate, applied to the frozen manifest, yields pure-plus-majority supervised totals 6,955/1,751/3,572. The excess is numerically equal to ambiguous counts 904/205/479 for the respective partitions. That relation is observed, but available provenance does not establish the historical generation rule or explain why the values differ. This study retains the discrepancy as OPEN instead of assigning an unverified cause.

For reproducibility, the audit source is \code{scripts/week2_data_audit.py}; its machine-readable receipt is \code{output/week2_foundations_20261103/data_audit.json}. Per-pair aligned overlap evidence is in \code{data_aligned_overlap_pairs.csv}; the support snapshot comparison is in \code{data_support_comparison.json}. The run recorded source and input hashes, timestamp, and environment. The source revision identified for this study is \code{633dc05a80511af4d4580cd985edd494cdc373ff} on \code{main}; that identifies the frozen baseline used by the audit, while this report edit is a later documentation change. Re-running the read-only audit can establish reproducibility against files present at rerun time, but cannot retroactively resolve missing upstream provenance.

\section{W2-08: Final data contract and handoff}
The following contract is the data-side handoff to the companion CNN study. It records shape, ordering, units, identity, label eligibility, normalization, and stream semantics. The entry marked specified is a requirement for an integration boundary; it is not evidence that a runtime implementation already enforces it.

{\small\renewcommand{\arraystretch}{1.08}
\begin{longtable}{@{}P{2.65cm}P{6.2cm}P{6.1cm}@{}}\toprule
Field & Frozen data contract & Evidence state / handoff rule\\\midrule\endhead
Tensor shape and type & Stored raw windows $X$ and normalized windows $Z=N(X)$ each use float32 with shape $[6,128]$; batches are $[B,6,128]$. The CNN consumes $Z$, after applying the frozen transform exactly once. Values must be finite. & Shape, dtype, finiteness, and exact reconstruction VERIFIED for the audited frozen artifact. Reject malformed shape or nonfinite values at the future boundary; enforcement NOT IMPLEMENTED here.\\
Channel order & \code{acc_x, acc_y, acc_z, gyro_x, gyro_y, gyro_z}. & Source order VERIFIED. Preserve order when forming a model tensor.\\
Rate and physical units & 50 Hz; acceleration is stored in multiples of $g$; gyroscope angular velocity in rad/s. & Local source contract VERIFIED; no unit conversion is performed in the audited path.\\
Window rule & Complete windows of 128 samples, starts every 64 samples, independently within each recording; no cross-recording windows. & Source rule and artifact reconstruction VERIFIED. Window duration is 2.56 s by sample slots, endpoint separation 2.54 s; stride duration is 1.28 s.\\
Recording identity & Every window remains associated with its recording and subject identifiers and split assignment. & Manifest/split reconstruction VERIFIED. A future streaming integration must reset window accumulation at recording boundary or explicit stream reset; runtime reset handling NOT IMPLEMENTED here.\\
Label policy & Class IDs 1--12 in the locked order; pure windows and accepted majority windows with at least 77 of 128 samples assigned to the dominant class are supervised. Ambiguous and unlabeled rows remain in the stream but have no single accepted target. & Current local policy VERIFIED. Upstream annotation index origin and historical support provenance remain OPEN.\\
Normalization & Statistics are computed only from 6,955 supervised training windows over window and time axes; float64 accumulation, float32 stored values, standard-deviation floor $10^{-6}$, then float32 output. Preprocessing ID: \code{aa8e1cc20829307187e0a78200f9957f7dd068b1dc26063c2720ad832bb2fd6e}. & Recomputed means, standard deviations, and ID VERIFIED against saved normalization and selected checkpoint metadata. Validation/test-specific fitting is excluded.\\
Denominators & Full stream counts are train/validation/test 9,978/2,422/5,053. Supervised counts are 6,955/1,751/3,572; ambiguous and unlabeled counts remain separately recorded. & Counts VERIFIED against the manifest and current source predicate. Preserve both stream and supervised denominators in downstream analysis.\\
Provenance and quality & Bind downstream artifacts to source revision, input hashes, builder, manifest, normalization ID, and audit receipt. Preserve the G2 transition-support HOLD and the open annotation/support provenance questions. & Existing receipts VERIFIED as file evidence. Formal annotation validity, representative-population status, and user acceptance are NOT ESTABLISHED.\\
\bottomrule\end{longtable}}

This handoff names recording-boundary reset and malformed-input rejection as required stream semantics, because windows must not cross recordings and downstream consumers need a well-formed tensor. Those behaviors are specification only in this document; this audit does not exercise or claim an implemented runtime reset, mask, fallback, or error path. Any mask supplied by an integration layer must retain its denominator and policy meaning rather than silently changing support counts. The data contract also does not define caching or certificate lifecycle; those are responsibilities of the separate system/model work.

\section{Conclusions and evidence status}
The study objective was to establish what data enter the frozen reference and whether the local source-to-artifact path can be reproduced. The reconstruction audit verifies the physical six-channel stream, 128/64 window coordinates, current label-policy behavior, subject-disjoint split counts, supervised support, and fixed train-only normalization. It also verifies that aligned shared samples match exactly across all 17,392 eligible adjacent pairs. The data tensor is therefore specified with a concrete shape, order, rate, units, recording identity, target eligibility, and preprocessing identity for handoff.

The evidence does not settle whether the local annotation coordinate origin matches the upstream intended semantics; an offset changes 98 assignments. It does not explain the historical support snapshot discrepancy, certify transition data quality, or establish that the subjects represent a target deployment population. Overlapping windows and repeated subject measurements remain dependencies that must be reflected in any later statistical design. These are material limits on interpretation, even though local reconstruction itself is exact.

The window-count formula, shifted-overlap identity, and the preservation of shared values by a fixed channelwise affine transform are DERIVED under the stated indexing and fixed-statistic assumptions. Their numerical realization is separately VERIFIED on the frozen corpus. A derivation over real coordinates and an executed array-equality check are different evidence.

The completion ledger distinguishes verified local results from open or merely specified statements. CREATED means the report and receipts exist; VERIFIED means a named local check supports the claim; NOT ESTABLISHED means the available evidence cannot support it. No item is labeled ACCEPTED: acceptance requires review by an authorized person. The companion CNN report independently owns model architecture, operation counts, prediction reconstruction, accuracy and classwise performance, and margin analysis. This data report owns the input data contract and does not reproduce those model results.

\begin{longtable}{@{}P{1.5cm}P{3.3cm}P{6.8cm}P{3.3cm}@{}}\toprule
Unit & Data-study question & Result and evidence & Status\\\midrule\endhead
W2-01 & What physical data exist? & 30 subjects, 61 recordings, 122 sensor files, six inertial channels at 50 Hz; 561 engineered features are distinct from the six-channel input. & VERIFIED local inventory\\
W2-02 & How are windows formed? & 128 samples, stride 64; 17,453 float32 windows; 17,392 aligned pairs; exact rebuild and no boundary crossing. & VERIFIED reconstruction\\
W2-03 & How are labels attached? & Pure, majority, ambiguous, and unlabeled policy; majority threshold 77/128; one-sample alternative changes 98 assignments. & Policy VERIFIED; origin OPEN\\
W2-04 & How are data partitioned and counted? & Subject-disjoint 17/4/9 groups; stream and supervised counts reconcile; complete class support retained. & Counts VERIFIED; historical provenance OPEN; G2 HOLD\\
W2-05 & How is scale set? & Train-only means/standard deviations, precision, floor, and full preprocessing ID reproduced exactly. & VERIFIED against frozen records\\
W2-06 & What quality claims do dependencies allow? & Exact overlap, channelwise differences, lag summaries, and sparse transition support documented without independence or population claims. & Audit VERIFIED; inference bounded\\
W2-07 & What provenance is reproducible? & Local tensor/manifest reconstruction and support comparison recorded with hashes; annotation origin and old support generation unresolved. & Local results VERIFIED; provenance OPEN\\
W2-08 & What must the handoff preserve? & Shape/type, order, units/rate, window/reset boundary, label policy, normalization, denominators, and provenance contract specified. & Contract specified; runtime enforcement NOT IMPLEMENTED\\
\bottomrule\end{longtable}

From the repository root, the read-only data audit is reproduced with \code{.venv/bin/python scripts/week2_data_audit.py}. Its JSON records input hashes and runtime details. This command does not train or modify the model. The ledger reports data-study completion only; it makes no claim about CNN analysis or system-level acceptance.

\end{document}
